\documentclass[12pt, a4paper]{article}
\usepackage[T2A]{fontenc}
\usepackage[utf8]{inputenc}
\usepackage[russian]{babel}
\usepackage{amsmath, amssymb} % Для математических символов
\usepackage{enumitem} % Для гибкой настройки списков
\usepackage{geometry} % Для настройки полей
\geometry{top=2cm, bottom=2cm, left=2.5cm, right=2cm}

\begin{document}

\section*{Георгий Игоревич Вольфсон <<Делимость с человеческим лицом>>}

\subsection*{Задачи вузовских олимпиад}

\subsubsection*{Задача 1}
\textbf{Найдите сумму всех натуральных чисел, не превосходящих 2019, которые могут быть представлены в виде суммы двух взаимно простых чисел, больших 1.}

\textbf{Решение:}
Докажем, что все натуральные числа, кроме 1, 2, 3, 4 и 6, представимы в виде суммы двух взаимно простых чисел, больших 1.

\begin{itemize}
    \item \textbf{Нечётные $n \ge 5$:} $n = 2 + (n-2)$. Числа 2 и $n-2$ взаимно просты (одно чётное, другое нечётное), оба больше 1.
    \item \textbf{Чётные $n \ge 8$, не кратные 3:} $n = 3 + (n-3)$. Числа 3 и $n-3$ взаимно просты (если бы общий делитель был больше 1, то это было бы 3, но $n-3$ не делится на 3). Оба больше 1.
    \item \textbf{Чётные $n \ge 10$, кратные 3, но не кратные 5:} $n = 5 + (n-5)$. Числа 5 и $n-5$ взаимно просты (аналогично). Оба больше 1.
    \item \textbf{Чётные $n \ge 12$:} по постулату Бертрана существует простое число $p$ с $n/2 < p < n$. Тогда $p \nmid n$, и $n = p + (n-p)$, где $\text{НОД}(p, n-p) = 1$. Оба числа больше 1.
\end{itemize}

Непосредственной проверкой убеждаемся, что числа 1, 2, 3, 4, 6 не представимы, а 5, 7, 8, 9, 10, 11, 12 представимы. Значит, единственные исключения: 1, 2, 3, 4, 6.

Сумма всех чисел от 1 до 2019:
\[
\frac{2019 \cdot 2020}{2} = 2\,039\,190.
\]
Вычитаем сумму исключений:
\[
1+2+3+4+6 = 16.
\]
\[
2\,039\,190 - 16 = 2\,039\,174.
\]

\textbf{Ответ:} 2\,039\,174.

\subsubsection*{Задача 2}
\textbf{Сколько существует различных пар целых чисел $(x,y)$, являющихся делителями числа 540, для которых НОД$(x, y) = 2$? Пары $(x,y)$ и $(y,x)$ считайте одной парой.}

\textbf{Решение:}
$540 = 2^2 \cdot 3^3 \cdot 5$. Пусть $x = 2^a \cdot 3^b \cdot 5^c$, $y = 2^d \cdot 3^e \cdot 5^f$, где $a,d \ge 1$, $\min(a,d) = 1$ (степень двойки в НОД ровно 1), $\min(b,e) = 0$, $\min(c,f) = 0$. Также $a,d \le 2$, $b,e \le 3$, $c,f \le 1$.

Случай $a = d = 1$: для 3 — 7 вариантов, для 5 — 3 варианта, итого 21 упорядоченная пара.
Случай $a = 1, d = 2$ или $a = 2, d = 1$: для 3 — 7 вариантов, для 5 — 3 варианта, итого $2 \cdot 21 = 42$ упорядоченных пары.
Всего упорядоченных пар: $21 + 42 = 63$. Из них пара $(2,2)$ — единственная с $x = y$. Остальные 62 разбиваются на пары $(x,y)$ и $(y,x)$. Неупорядоченных пар: $62/2 + 1 = 32$.

\textbf{Ответ:} 32.

\subsubsection*{Задача 3}
\textbf{Натуральное число $a$ раскладывается в произведение трёх различных простых делителей. Сумма всех его делителей, включая 1 и $a$, равна 72. Найдите число $a$.}

\textbf{Решение:}
Пусть $a = xyz$, где $x, y, z$ — различные простые числа. Сумма всех делителей равна:
\[
(1+x)(1+y)(1+z) = 72 = 2^3 \cdot 3^2.
\]
Перебираем разложения 72 на три множителя, каждый из которых на 1 больше простого числа:
\begin{itemize}
    \item $2 \cdot 4 \cdot 9$: $9 = 1+8$, 8 не простое.
    \item $3 \cdot 4 \cdot 6$: $3 = 1+2$, $4 = 1+3$, $6 = 1+5$ — все простые.
\end{itemize}
Значит, простые делители: 2, 3, 5. Число $a = 2 \cdot 3 \cdot 5 = 30$.

\textbf{Ответ:} 30.

\subsubsection*{Задача 4}
\textbf{Сколько различных пар целых чисел $1 \le x, y \le 100$ удовлетворяют уравнению НОД$(\log_2 x; 36) = (\log_3 y)^2$? Найдите эти пары.}

\textbf{Решение:}
$\log_2 x$ принимает целые значения: 0, 1, 2, 3, 4, 5, 6 (для $x = 1, 2, 4, 8, 16, 32, 64$).

Проверяем каждый случай:
\begin{itemize}
    \item $x=1$: $\log_2 x = 0$, НОД$(0,36)=36$, $(\log_3 y)^2 = 36 \Rightarrow \log_3 y = \pm 6$. $y = 3^6 = 729 > 100$ — не подходит.
    \item $x=2$: $\log_2 x = 1$, НОД$(1,36)=1$, $(\log_3 y)^2 = 1 \Rightarrow \log_3 y = \pm 1$. $y = 3$ подходит. Пара $(2;3)$.
    \item $x=4$: $\log_2 x = 2$, НОД$(2,36)=2$, $(\log_3 y)^2 = 2 \Rightarrow \log_3 y = \pm\sqrt{2}$ — не целое.
    \item $x=8$: $\log_2 x = 3$, НОД$(3,36)=3$, $(\log_3 y)^2 = 3 \Rightarrow \log_3 y = \pm\sqrt{3}$ — не целое.
    \item $x=16$: $\log_2 x = 4$, НОД$(4,36)=4$, $(\log_3 y)^2 = 4 \Rightarrow \log_3 y = \pm 2$. $y = 9$ подходит. Пара $(16;9)$.
    \item $x=32$: $\log_2 x = 5$, НОД$(5,36)=1$, $(\log_3 y)^2 = 1 \Rightarrow y = 3$. Пара $(32;3)$.
    \item $x=64$: $\log_2 x = 6$, НОД$(6,36)=6$, $(\log_3 y)^2 = 6 \Rightarrow \log_3 y = \pm\sqrt{6}$ — не целое.
\end{itemize}

\textbf{Ответ:} 3 пары: $(2;3)$, $(16;9)$, $(32;3)$.

\end{document}